% Part: first-order-logic
% Chapter: beyond
% Section: modal-logics

\documentclass[../../../include/open-logic-section]{subfiles}

\begin{document}

\olfileid{fol}{byd}{mod}

\olsection{Logika Modal}

Gatekna tuladha ukara kondisional iki:
\begin{quote}
  Yen Jeremy dhewekan ing kamar iku,
  mula dheweke mabuk lan wuda lan
  njogèd ing kursi-kursi.
\end{quote}
Iki tuladha pratelan kondisional kang
mbokmenawa bener sacara material,
nanging isih bisa nyasarake.
Ukara kasebut kaya-kaya nuduhake
sesambungan kang luwih kuwat
antarane prasyarat lan akibat
tinimbang mung kahanan yen
prasyarate luput utawa akibate
bener. Tegese, susunan ukara
iku nuduhake yen pratelan
kasebut ora mung bener ing
donya iki (ing ngendi bisa
bener kanthi sepele amarga
Jeremy ora dhewekan ing kamar),
nanging akibate \emph{mesthine}
bener \emph{yen} prasyarate
bener. Kanthi tembung liya,
pratelan iku bisa ditegesi
bener ora mung ing donya iki,
nanging ing saben donya kang
"mungkin"; utawa pratelan
iku \emph{mesthi} bener,
dudu mung bener ing donya
tartamtu iki.

Logika modal dirancang kanggo
njlentrehake keniscayan kaya
mangkono. Logika proposisional
modal dioleh saka logika
proposisional lumrah kanthi
nambah operator kothak:
yen $!A$ iku !!a{formula},
mula $\Box !A$ uga
formula. Miturut
teges kang dikarepake,
$\Box !A$ nyatakake
yen $!A$ \emph{mesthi}
bener, utawa bener ing
saben donya kang mungkin.
$\Diamond !A$ lumrahe
dianggep cekakan kanggo
$\lnot \Box \lnot !A$,
lan bisa diwaca minangka
pratelan yen $!A$
\emph{bisa uga} bener.
Mesthi wae, modalitas
uga bisa ditambahake
marang logika predikat.

!!{structure}s Kripke bisa
dienggo kanggo menehi
semantik logika modal;
malah Kripke wiwitane
ngrancang semantik iki
kanthi logika modal
ing pikirane. Tinimbang
mbatesi mung urutan
parsial, sacara luwih
umum ana himpunan
"donya-donya kang mungkin",
$P$, lan relasi binari
"aksesibilitas"
$\Atom{R}{x,y}$
antarane donya.
Miturut pangertèn
lumrah, $\Atom{R}{p,q}$
nyatakake yen donya~$q$
cocog karo~$p$;
yaiku, yen kita "ana"
ing donya~$p$, kita
kudu nganggep mungkin
yen donya iku bisa
kaya~$q$.

Logika modal sok diarani
logika "intensional",
minangka kosok baline
logika "ekstensional".
Semantik kang dimaksud
kanggo logika ekstensional,
kaya logika klasik,
mung ngacu marang siji
donya, yaiku donya
"aktual"; dene semantik
logika intensional
gumantung marang ontologi
kang luwih rumit.
Kejaba kanggo nyusun
!!{structure} keniscayan,
modalitas uga bisa
kanggo nyusun
!!{structure} wangun
basa liyane, kanthi
nafsirake maneh
$\Box$ lan $\Diamond$
manut panganggone.
Upamane:
\begin{enumerate}
\item Ing logika keterbuktian,
  $\Box !A$ diwaca
  "$!A$ bisa dibuktekake"
  lan $\Diamond !A$ diwaca
  "$!A$ konsisten."
\item Ing logika epistemik,
  $\Box !A$ bisa diwaca
  "Aku ngerti $!A$"
  utawa "Aku percaya $!A$."
\item Ing logika temporal,
  $\Box !A$ bisa diwaca
  "$!A$ tansah bener"
  lan $\Diamond !A$ diwaca
  "$!A$ kadhang bener."
\end{enumerate}

Kita kepengin nambah aturan
lan aksioma kang ngurusi
modalitas. Upamane, sistem
$\Log{S4}$ dumadi saka
aksioma lan aturan lumrah
logika proposisional,
bebarengan karo aksioma
ing ngisor iki:
\begin{align*}
& \Box (!A \lif !B) \lif (\Box !A \lif \Box !B)\\
& \Box !A \lif !A \\
& \Box !A \lif \Box \Box !A
\intertext{sarta aturan, "saka teorema $!A$ kang dibuktekake tanpa asumsi mbukak, simpulake $\Box !A$."
  $\Log{S5}$ nambah aksioma ing ngisor iki:}
& \Diamond !A \lif \Box \Diamond !A
\end{align*}
Cathetan cakupan SOL6-F016: aturan necessitation iki mung ditrapake marang teorema tanpa asumsi mbukak, ora marang pratelan kang mung dadi premis. Yen saka premis sembarang diidinake, formula bener ing donya saiki bakal dipaksa bener ing kabeh donya kang bisa diakses, kang ora sah.
Wujud liya saka aksioma
iki bisa cocog kanggo
panganggone kang beda-beda;
upamane, S5 lumrahe
dianggep nggambarake
keniscayan logis.
Kang migunani, kita
lumrahe bisa nemokake
semantik kang ndadekake
sistem !!{derivation}
iku sah lan jangkep
kanthi mbatesi relasi
aksesibilitas ing
!!{structure}s Kripke
kanthi cara lumrah.
Upamane, $\Log{S4}$
cocog karo kelas
!!{structure}s Kripke
kang relasi aksesibilitase
refleksif lan transitif.
$\Log{S5}$ cocog
karo kelas
!!{structure}s Kripke
kang relasi aksesibilitase
{\em universal}, yaiku
saben donya bisa diakses
saka saben donya liyane;
mula $\Box !A$ bener
yen lan mung yen $!A$
bener ing saben donya.

\end{document}
